\documentclass[aspectratio=169]{beamer}\usepackage[]{graphicx}\usepackage[]{xcolor}
% maxwidth is the original width if it is less than linewidth
% otherwise use linewidth (to make sure the graphics do not exceed the margin)
\makeatletter
\def\maxwidth{ %
  \ifdim\Gin@nat@width>\linewidth
    \linewidth
  \else
    \Gin@nat@width
  \fi
}
\makeatother

\definecolor{fgcolor}{rgb}{0.345, 0.345, 0.345}
\newcommand{\hlnum}[1]{\textcolor[rgb]{0.686,0.059,0.569}{#1}}%
\newcommand{\hlsng}[1]{\textcolor[rgb]{0.192,0.494,0.8}{#1}}%
\newcommand{\hlcom}[1]{\textcolor[rgb]{0.678,0.584,0.686}{\textit{#1}}}%
\newcommand{\hlopt}[1]{\textcolor[rgb]{0,0,0}{#1}}%
\newcommand{\hldef}[1]{\textcolor[rgb]{0.345,0.345,0.345}{#1}}%
\newcommand{\hlkwa}[1]{\textcolor[rgb]{0.161,0.373,0.58}{\textbf{#1}}}%
\newcommand{\hlkwb}[1]{\textcolor[rgb]{0.69,0.353,0.396}{#1}}%
\newcommand{\hlkwc}[1]{\textcolor[rgb]{0.333,0.667,0.333}{#1}}%
\newcommand{\hlkwd}[1]{\textcolor[rgb]{0.737,0.353,0.396}{\textbf{#1}}}%
\let\hlipl\hlkwb

\usepackage{framed}
\makeatletter
\newenvironment{kframe}{%
 \def\at@end@of@kframe{}%
 \ifinner\ifhmode%
  \def\at@end@of@kframe{\end{minipage}}%
  \begin{minipage}{\columnwidth}%
 \fi\fi%
 \def\FrameCommand##1{\hskip\@totalleftmargin \hskip-\fboxsep
 \colorbox{shadecolor}{##1}\hskip-\fboxsep
     % There is no \\@totalrightmargin, so:
     \hskip-\linewidth \hskip-\@totalleftmargin \hskip\columnwidth}%
 \MakeFramed {\advance\hsize-\width
   \@totalleftmargin\z@ \linewidth\hsize
   \@setminipage}}%
 {\par\unskip\endMakeFramed%
 \at@end@of@kframe}
\makeatother

\definecolor{shadecolor}{rgb}{.97, .97, .97}
\definecolor{messagecolor}{rgb}{0, 0, 0}
\definecolor{warningcolor}{rgb}{1, 0, 1}
\definecolor{errorcolor}{rgb}{1, 0, 0}
\newenvironment{knitrout}{}{} % an empty environment to be redefined in TeX

\usepackage{alltt}




\usetheme{default}
% Slide setup, colour independent

\usepackage{amsmath,amssymb,amsthm}
\usepackage[utf8]{inputenc}
\usepackage{colortbl}
\usepackage{bm}
\usepackage{xcolor}
\usepackage{dsfont}
\usepackage{setspace}
% To use \ding{234} and the like
\usepackage{pifont}
% To cross reference between slide files
\usepackage{zref-xr,zref-user}
% Use something like
% \zexternaldocument{fileI}
% in the tex files. And cite using \zref instead of \ref

% Cross-reference system - see CROSS-REFERENCE-SETUP.md for manual setup instructions
\usepackage{booktabs}
\usepackage{marvosym}
\usepackage{cancel}
%\usepackage{transparent}
% Make doi clickable in the bibliography?
\usepackage{doi}

\usepackage[T1]{fontenc}

\usepackage{longtable}

% For heavier titles
\usepackage{helvet} % Enables Helvetica font family


% Fields and the like
\def\IC{\mathbb{C}}
\def\IE{\mathbb{E}}
\def\IF{\mathbb{F}}
\def\II{\mathbb{I}}
\def\IJ{\mathbb{J}}
\def\IK{\mathbb{K}}
\def\IM{\mathbb{M}}
\def\IN{\mathbb{N}}
\def\IP{\mathbb{P}}
\def\IR{\mathbb{R}}
\newcommand{\IRplus}{\mathbb{R}_{\ge 0}}
\def\IZ{\mathbb{Z}}
\def\11{\mathds{1}}


% Bold lowercase
\def\ba{\bm{a}}
\def\bb{\bm{b}}
\def\bc{\bm{c}}
\def\bd{\bm{d}}
\def\be{\bm{e}}
\def\bf{\bm{f}}
\def\bg{\bm{g}}
\def\bh{\bm{h}}
\def\bi{\bm{i}}
\def\bj{\bm{j}}
\def\bk{\bm{k}}
\def\bn{\bm{n}}
\def\bp{\bm{p}}
\def\br{\bm{r}}
\def\bs{\bm{s}}
\def\bu{\bm{u}}
\def\bv{\bm{v}}
\def\bw{\bm{w}}
\def\bx{\bm{x}}
\def\by{\bm{y}}
\def\bz{\bm{z}}
\newcommand{\vect}[1]{\bm{#1}}

% Bold capitals
\def\bB{\bm{B}}
\def\bD{\bm{D}}
\def\bE{\bm{E}}
\def\bF{\bm{F}}
\def\bG{\bm{G}}
\def\bI{\bm{I}}
\def\bL{\bm{L}}
\def\bN{\bm{N}}
\def\bP{\bm{P}}
\def\bR{\bm{R}}
\def\bS{\bm{S}}
\def\bT{\bm{T}}
\def\bX{\bm{X}}

% Bold numbers
\def\b0{\bm{0}}

% Bold greek
\bmdefine{\bmu}{\bm{\mu}}
\def\bphi{\bm{\phi}}
\def\bvarphi{\bm{\varphi}}
\def\bPi{\bm{\Pi}}
\def\bGamma{\bm{\Gamma}}

% Bold red sentence
\def\boldred#1{{\color{red}\textbf{#1}}}
\def\defword#1{{\color{orange}\textbf{#1}}}

% Caligraphic letters
\def\A{\mathcal{A}}
\def\B{\mathcal{B}}
\def\C{\mathcal{C}}
\def\D{\mathcal{D}}
\def\E{\mathcal{E}}
\def\F{\mathcal{F}}
\def\G{\mathcal{G}}
\def\H{\mathcal{H}}
\def\I{\mathcal{I}}
\def\L{\mathcal{L}}
\def\M{\mathcal{M}}
\def\N{\mathcal{N}}
\def\P{\mathcal{P}}
\def\R{\mathcal{R}}
\def\S{\mathcal{S}}
\def\T{\mathcal{T}}
\def\U{\mathcal{U}}
\def\V{\mathcal{V}}

% Adding space for prime (') where needed
\def\pprime{\,'}
% Adding space for star (\star) where needed
\def\pstar{{\,\star}}

% tt font for code
\def\code#1{{\tt #1}}

% i.e., e.g.
\def\eg{\emph{e.g.}}
\def\ie{\emph{i.e.}}


% Operators and special symbols
\def\nbOne{{\mathchoice {\rm 1\mskip-4mu l} {\rm 1\mskip-4mu l}
{\rm 1\mskip-4.5mu l} {\rm 1\mskip-5mu l}}}
\def\cov{\ensuremath{\mathsf{cov}}}
\def\Var{\ensuremath{\mathsf{Var}\ }}
\def\Im{\textrm{Im}\;}
\def\Re{\textrm{Re}\;}
\def\det{\ensuremath{\mathsf{det}}}
\def\diag{\ensuremath{\mathsf{diag}}}
\def\nullspace{\ensuremath{\mathsf{null}}}
\def\nullity{\ensuremath{\mathsf{nullity}}}
\def\rank{\ensuremath{\mathsf{rank}}}
\def\range{\ensuremath{\mathsf{range}}}
\def\sgn{\ensuremath{\mathsf{sgn}}}
\def\Span{\ensuremath{\mathsf{span}}}
\def\tr{\ensuremath{\mathsf{tr}}}
\def\imply{$\Rightarrow$}
\def\restrictTo#1#2{\left.#1\right|_{#2}}
\newcommand{\parallelsum}{\mathbin{\!/\mkern-5mu/\!}}
\def\dsum{\mathop{\displaystyle \sum }}%
\def\dind#1#2{_{\substack{#1\\ #2}}}

\newcommand{\Qmatrix}[1]{%
  \begin{pmatrix}#1\end{pmatrix}%
}

\DeclareMathOperator{\GL}{GL}
\DeclareMathOperator{\Rel}{Re}
\def\Nt#1{\left|\!\left|\!\left|#1\right|\!\right|\!\right|}
\newcommand{\tripbar}{|\! |\! |}



% The beamer bullet (in base colour)
\def\bbullet{\leavevmode\usebeamertemplate{itemize item}\ }

% Theorems and the like
\newtheorem{proposition}[theorem]{Proposition}
\newtheorem{property}[theorem]{Property}
\newtheorem{importantproperty}[theorem]{Property}
\newtheorem{importanttheorem}[theorem]{Theorem}
%\newtheorem{lemma}[theorem]{Lemma}
%\newtheorem{corollary}[theorem]{Corollary}
\newtheorem{remark}[theorem]{Remark}
\setbeamertemplate{theorems}[numbered]
%\setbeamertemplate{theorems}[ams style]

%
%\usecolortheme{orchid}
%\usecolortheme{orchid}

\def\red{\color[rgb]{1,0,0}}
\def\blue{\color[rgb]{0,0,1}}
\def\green{\color[rgb]{0,1,0}}

% Fix skipping lines after items in the bibliography
\setbeamertemplate{bibliography entry title}{}
\setbeamertemplate{bibliography entry location}{}
\setbeamertemplate{bibliography entry note}{}

% Get rid of navigation stuff
\setbeamertemplate{navigation symbols}{}

% Set footline/header line
\setbeamertemplate{footline}
{%
\quad p. \insertpagenumber \quad--\quad \insertsection\vskip2pt
}
% \setbeamertemplate{headline}
% {%
% \quad\insertsection\hfill p. \insertpagenumber\quad\mbox{}\vskip2pt
% }


\makeatletter
\newlength\beamerleftmargin
\setlength\beamerleftmargin{\Gm@lmargin}
\makeatother

% Colours for special pages
\def\extraContent{yellow!20}


%%%%%%%%%%%%%%%%%
\usepackage{tikz}
\usetikzlibrary{shapes,arrows}
\usetikzlibrary{positioning}
\usetikzlibrary{shapes.symbols,shapes.callouts,patterns}
\usetikzlibrary{calc,fit}
\usetikzlibrary{backgrounds}
\usetikzlibrary{decorations.pathmorphing,fit,petri}
\usetikzlibrary{automata}
\usetikzlibrary{fadings}
\usetikzlibrary{patterns,hobby}
\usetikzlibrary{backgrounds,fit,petri}
\usetikzlibrary{tikzmark}

\usepackage{pgfplots}
\pgfplotsset{compat=1.6}
\pgfplotsset{ticks=none}

\usetikzlibrary{decorations.markings}
\usetikzlibrary{arrows.meta}
\tikzset{>=stealth}

% For tikz
\tikzstyle{cloud} = [draw, ellipse,fill=red!20, node distance=0.87cm,
minimum height=2em]
\tikzstyle{line} = [draw, -latex']


%%% For max frame images
\newenvironment{changemargin}[2]{%
\begin{list}{}{%
\setlength{\topsep}{0pt}%
\setlength{\leftmargin}{#1}%
\setlength{\rightmargin}{#2}%
\setlength{\listparindent}{\parindent}%
\setlength{\itemindent}{\parindent}%
\setlength{\parsep}{\parskip}%
}%
\item[]}{\end{list}}


% Make one image take up the entire slide content area in beamer,.:
% centered/centred full-screen image, with title:
% This uses the whole screen except for the 1cm border around it
% all. 128x96mm
\newcommand{\titledFrameImage}[2]{
\begin{frame}{#1}
%\begin{changemargin}{-1cm}{-1cm}
\begin{center}
\includegraphics[width=108mm,height=\textheight,keepaspectratio]{#2}
\end{center}
%\end{changemargin}
\end{frame}
}

% Make one image take up the entire slide content area in beamer.:
% centered/centred full-screen image, no title:
% This uses the whole screen except for the 1cm border around it
% all. 128x96mm
\newcommand{\plainFrameImage}[1]{
\begin{frame}[plain]
%\begin{changemargin}{-1cm}{-1cm}
\begin{center}
\includegraphics[width=108mm,height=76mm,keepaspectratio]{#1}
\end{center}
%\end{changemargin}
\end{frame}
}

% Make one image take up the entire slide area, including borders, in beamer.:
% centered/centred full-screen image, no title:
% This uses the entire whole screen
\newcommand{\maxFrameImage}[1]{
\begin{frame}[plain]
\begin{changemargin}{-1cm}{-1cm}
\begin{center}
\includegraphics[width=\paperwidth,height=\paperheight,keepaspectratio]
{#1}
\end{center}
\end{changemargin}
\end{frame}
}

% This uses the entire whole screen (to include in frame)
\newcommand{\maxFrameImageNoFrame}[1]{
\begin{changemargin}{-1cm}{-1cm}
\begin{center}
\includegraphics[width=\paperwidth,height=0.99\paperheight,keepaspectratio]
{#1}
\end{center}
\end{changemargin}
}

% Make one image take up the entire slide area, including borders, in beamer.:
% centered/centred full-screen image, no title:
% This uses the entire whole screen
\newcommand{\maxFrameImageColor}[2]{
\begin{frame}[plain]
\setbeamercolor{normal text}{bg=#2!20}
\begin{changemargin}{-1cm}{-1cm}
\begin{center}
\includegraphics[width=\paperwidth,height=\paperheight,keepaspectratio]
{#1}
\end{center}
\end{changemargin}
\end{frame}
}


\usepackage{tikz}
\usetikzlibrary{patterns,hobby,matrix}
\usepackage{pgfplots}
\pgfplotsset{compat=1.6}
\pgfplotsset{ticks=none}

\usetikzlibrary{backgrounds}
\usetikzlibrary{decorations.markings}
\usetikzlibrary{arrows.meta}
\tikzset{>=stealth}

\tikzset{
  clockwise arrows/.style={
    postaction={
      decorate,
      decoration={
        markings,
        mark=between positions 0.1 and 0.9 step 40pt with {\arrow{>}},
   }}}}


% New integrated section command: creates section and section slide
\newcommand{\Ssection}[2]{
\section{#1}
\begin{frame}[noframenumbering,plain]
  \begin{tikzpicture}[remember picture,overlay]
    \node[above right,inner sep=0pt,opacity=0.2] at (current page.south west)
    {
        \includegraphics[height=\paperheight,width=\paperwidth]{#2}
    };
  \end{tikzpicture}
  \setbeamercolor{section in toc}{fg=section_page_list_colour}
  \setbeamerfont{section in toc}{size=\Large,series=\bfseries}
  \setbeamertemplate{section in toc shaded}[default][60]
  \tableofcontents[
    currentsection,
    sectionstyle=show/shaded,
    subsectionstyle=show/hide/hide,
    subsubsectionstyle=hide/hide/hide]
\end{frame}
\addtocounter{page}{-1}
}

% New integrated section command with subsections: creates section and section slide showing subsections
\newcommand{\SsectionWithSubs}[2]{
\section{#1}
\begin{frame}[noframenumbering,plain]
  \begin{tikzpicture}[remember picture,overlay]
    \node[above right,inner sep=0pt,opacity=0.2] at (current page.south west)
    {
        \includegraphics[height=\paperheight,width=\paperwidth]{#2}
    };
  \end{tikzpicture}
  \setbeamercolor{section in toc}{fg=section_page_list_colour}
  \setbeamerfont{section in toc}{size=\Large,series=\bfseries}
  \setbeamertemplate{section in toc shaded}[default][60]
  \tableofcontents[
    currentsection,
    sectionstyle=show/hide,
    subsectionstyle=show/show/hide,
    subsubsectionstyle=hide/hide/hide]
\end{frame}
\addtocounter{page}{-1}
}

% New integrated subsection command: creates subsection and subsection slide
\newcommand{\Ssubsection}[2]{
\subsection{#1}
\begin{frame}[noframenumbering,plain]
  \begin{tikzpicture}[remember picture,overlay]
    \node[above right,inner sep=0pt,opacity=0.2] at (current page.south west)
    {
        \includegraphics[height=\paperheight,width=\paperwidth]{#2}
    };
  \end{tikzpicture}
  \setbeamercolor{section in toc}{fg=subsection_page_list_colour}
  \setbeamerfont{section in toc}{size=\Large,series=\bfseries}
  \setbeamertemplate{section in toc shaded}[default][60]
  \setbeamerfont{subsection in toc}{series=\bfseries}
  \setbeamertemplate{subsection in toc shaded}[default][50]
  \tableofcontents[
    currentsection,
    sectionstyle=show/hide,
    subsectionstyle=show/shaded/hide,
    subsubsectionstyle=hide/hide/hide]
\end{frame}
\addtocounter{page}{-1}
}

% New integrated subsubsection command: creates subsubsection and subsubsection slide
\newcommand{\Ssubsubsection}[2]{
\subsubsection{#1}
\begin{frame}[noframenumbering,plain]
  \begin{tikzpicture}[remember picture,overlay]
    \node[above right,inner sep=0pt,opacity=0.2] at (current page.south west)
    {
        \includegraphics[height=\paperheight,width=\paperwidth]{#2}
    };
  \end{tikzpicture}
  \setbeamercolor{section in toc}{fg=subsub_header_section}
  \setbeamerfont{section in toc}{size=\Large,series=\bfseries}
  \setbeamertemplate{section in toc shaded}[default][60]
  \setbeamerfont{subsection in toc}{series=\bfseries}
  \setbeamertemplate{subsection in toc shaded}[default][50]
  \setbeamertemplate{subsubsection in toc shaded}[default][50]
  \tableofcontents[
    currentsection,
    sectionstyle=show/hide,
    subsectionstyle=show/hide/hide,
    subsubsectionstyle=show/shaded/hide]
\end{frame}
\addtocounter{page}{-1}
}

% Legacy commands (kept for backward compatibility)
% Beginning of a section
\newcommand{\newSectionSlide}[1]{
\begin{frame}[noframenumbering,plain]
  \begin{tikzpicture}[remember picture,overlay]
    \node[above right,inner sep=0pt,opacity=0.2] at (current page.south west)
    {
        \includegraphics[height=\paperheight,width=\paperwidth]{#1}
    };
  \end{tikzpicture}
  \setbeamercolor{section in toc}{fg=section_page_list_colour}
  \setbeamerfont{section in toc}{size=\Large,series=\bfseries}
  \setbeamertemplate{section in toc shaded}[default][60]
  \tableofcontents[
    currentsection,
    sectionstyle=show/shaded,
    subsectionstyle=show/hide/hide,
    subsubsectionstyle=hide/hide/hide]
\end{frame}
\addtocounter{page}{-1}
}

% Beginning of a section in which we also show subsections
\newcommand{\newSectionWithSubsSlide}[1]{
	\begin{frame}[noframenumbering,plain]
		\begin{tikzpicture}[remember picture,overlay]
			\node[above right,inner sep=0pt,opacity=0.2] at (current page.south west)
			{
				\includegraphics[height=\paperheight,width=\paperwidth]{#1}
			};
		\end{tikzpicture}
		\setbeamercolor{section in toc}{fg=section_page_list_colour}
		\setbeamerfont{section in toc}{size=\Large,series=\bfseries}
		\setbeamertemplate{section in toc shaded}[default][60]
		\tableofcontents[
		currentsection,
		sectionstyle=show/hide,
		subsectionstyle=show/show/hide,
		subsubsectionstyle=hide/hide/hide]
	\end{frame}
	\addtocounter{page}{-1}
}

% Beginning of a subsection
\newcommand{\newSubSectionSlide}[1]{
\begin{frame}[noframenumbering,plain]
  \begin{tikzpicture}[remember picture,overlay]
    \node[above right,inner sep=0pt,opacity=0.2] at (current page.south west)
    {
        \includegraphics[height=\paperheight,width=\paperwidth]{#1}
    };
  \end{tikzpicture}
  \setbeamercolor{section in toc}{fg=subsection_page_list_colour}
  \setbeamerfont{section in toc}{size=\Large,series=\bfseries}
  \setbeamertemplate{section in toc shaded}[default][60]
  \setbeamerfont{subsection in toc}{series=\bfseries}
  \setbeamertemplate{subsection in toc shaded}[default][50]
  \tableofcontents[
    currentsection,
    sectionstyle=show/hide,
    subsectionstyle=show/shaded/hide,
    subsubsectionstyle=hide/hide/hide]
\end{frame}
\addtocounter{page}{-1}
}

% Beginning of a subsubsection
\newcommand{\newSubSubSectionSlide}[1]{
\begin{frame}[noframenumbering,plain]
  \begin{tikzpicture}[remember picture,overlay]
    \node[above right,inner sep=0pt,opacity=0.2] at (current page.south west)
    {
        \includegraphics[height=\paperheight,width=\paperwidth]{#1}
    };
  \end{tikzpicture}
  \setbeamercolor{section in toc}{fg=subsub_header_section}
  \setbeamerfont{section in toc}{size=\Large,series=\bfseries}
  \setbeamertemplate{section in toc shaded}[default][60]
  \setbeamerfont{subsection in toc}{series=\bfseries}
  \setbeamertemplate{subsection in toc shaded}[default][50]
  \setbeamertemplate{subsubsection in toc shaded}[default][50]
  \tableofcontents[
    currentsection,
    sectionstyle=show/hide,
    subsectionstyle=show/hide/hide,
    subsubsectionstyle=show/shaded/hide]
\end{frame}
\addtocounter{page}{-1}
}


   %%%%%%%%%%%
% To have links to parts in the outline
\makeatletter
\AtBeginPart{%
  \addtocontents{toc}{\protect\beamer@partintoc{\the\c@part}{\beamer@partnameshort}{\the\c@page}}%
}
%% number, shortname, page.
\providecommand\beamer@partintoc[3]{%
  \ifnum\c@tocdepth=-1\relax
    % requesting onlyparts.
    \makebox[6em]{Part #1:} \textcolor{green!30!blue}{\hyperlink{#2}{#2}}
    \par
  \fi
}
\define@key{beamertoc}{onlyparts}[]{%
  \c@tocdepth=-1\relax
}
\makeatother%

\newcommand{\nameofthepart}{}
\newcommand{\nupart}[1]%
    {   \part{#1}%
        \renewcommand{\nameofthepart}{#1}%
        {
          \setbeamercolor{background canvas}{bg=orange!50}
          \begin{frame}{#1}%\partpage 
          \hypertarget{\nameofthepart}{}\tableofcontents%
          \end{frame}
        }
    }

% This command creates a title page using TikZ only
\newcommand{\tikztitlepage}[1]{%
\begin{frame}[plain,noframenumbering]
  \begin{tikzpicture}[remember picture,overlay]
    % Background image
    \node[above right,inner sep=0pt,opacity=0.1] 
      at (current page.south west) 
      {\includegraphics[width=\paperwidth,height=\paperheight]{#1}};

    % University logo
    \node[anchor=north east, inner sep=5pt, opacity=0.9] 
      at (current page.north east)
      {\includegraphics[width=0.2\textwidth]{FIGS-slides-admin/UM-logo-horizontal-CMYK.png}};
    
    % Title
    \node[anchor=center, align=center, 
          font=\fontsize{13}{15}\bfseries\color{UMbrown}, 
          text width=0.9\textwidth] 
          at ([yshift=2cm]current page.center)
          {\inserttitle};

      % Authors
      \node[anchor=center, align=center,
        font=\fontsize{10}{12}\bfseries\color{UMbrown},
        text width=0.7\textwidth]
        at ([yshift=0.8cm]current page.center)
        {\insertauthor};

      % Affiliation
      \node[anchor=north, align=center,
        font=\fontsize{9}{11}\color{UMbrown},
        text width=0.7\textwidth]
        at ([yshift=-0.2cm]current page.center)
        {\insertaffiliation};      
    % Date
    \node[anchor=north, align=center, 
          font=\fontsize{12}{16}\bfseries\color{UMbrown},
          text width=0.7\textwidth] 
          at ([yshift=0.2cm]current page.center)
          {\insertdate};

    % Land acknowledgement
    \node[anchor=south, align=justify, 
          font=\footnotesize, text=black, 
          text width=1.1\textwidth] 
          at ([yshift=0.5cm]current page.south)
          {The University of Manitoba campuses are located on original lands of Anishinaabeg, Ininew, Anisininew, Dakota and Dene peoples, and on the National Homeland of the Red River Métis.\\
          We respect the Treaties that were made on these territories, we acknowledge the harms and mistakes of the past, and we dedicate ourselves to move forward in partnership with Indigenous communities in a spirit of Reconciliation and collaboration.};
  \end{tikzpicture}
  \addtocounter{page}{-1}
\end{frame}
}
% The title page with figure
% \newcommand{\titlepagewithfigure}[1]{%
%   \begin{frame}[noframenumbering,plain]
%     \begin{tikzpicture}[remember picture,overlay]
%       \node[above right,inner sep=0pt,opacity=0.1] at (current page.south west)
%       {
%           \includegraphics[height=\paperheight,width=\paperwidth]{#1}
%       };
%       \node[anchor=north east,
%       inner sep=5pt,
%       opacity=0.9] at (current page.north east)
%       {
%           \includegraphics[width=0.2\textwidth]{FIGS-slides-admin/UM-logo-horizontal-CMYK.png}
%       };
%       \node[anchor=south, 
%       align=justify, 
%       text=black, 
%       text width=1.1\textwidth,
%       font=\footnotesize]  (land_acknowledgement)
%       at (current page.south) 
%       {The University of Manitoba campuses are located on original lands of Anishinaabeg, Ininew, Anisininew, Dakota and Dene peoples, and on the National Homeland of the Red River Métis.
%       We respect the Treaties that were made on these territories, we acknowledge the harms and mistakes of the past, and we dedicate ourselves to move forward in partnership with Indigenous communities in a spirit of Reconciliation and collaboration.};  
%       % \node[align=center, anchor=south,
%       % above=0.5cm of land_acknowledgement,
%       % text=black,
%       % font=\bfseries] {\insertdate};
%   \end{tikzpicture}
%   \setbeamercolor{title}{fg=title_page_title_colour}
%   \setbeamerfont{title}{size=\Large,series=\bfseries}
%   \setbeamercolor{author}{fg=title_page_author_colour}
%   \setbeamerfont{author}{size=\large,series=\bfseries}
%   \setbeamercolor{institute}{fg=title_page_institute_colour}
%   \setbeamerfont{institute}{size=\large,series=\bfseries}
%   \setbeamercolor{date}{fg=title_page_date_colour}
%   \setbeamerfont{date}{series=\bfseries}
% 	\titlepage
% \end{frame}
% \addtocounter{page}{-1}
% }

\newcommand{\titlepagewithfigure}[1]{%
  \begin{frame}[noframenumbering,plain]
    \begin{tikzpicture}[remember picture,overlay]
      \node[above right,inner sep=0pt,opacity=0.1] at (current page.south west)
      {
          \includegraphics[height=\paperheight,width=\paperwidth]{#1}
      };
      \node[anchor=north east,
      inner sep=5pt,
      opacity=0.9] at (current page.north east)
      {
          \includegraphics[width=0.2\textwidth]{FIGS-slides-admin/UM-logo-horizontal-CMYK.png}
      };
      \node[anchor=south, 
      align=justify, 
      text=black, 
      text width=1.1\textwidth,
      font=\footnotesize]  (land_acknowledgement)
      at (current page.south) 
      {The University of Manitoba campuses are located on original lands of Anishinaabeg, Ininew, Anisininew, Dakota and Dene peoples, and on the National Homeland of the Red River Métis.
      We respect the Treaties that were made on these territories, we acknowledge the harms and mistakes of the past, and we dedicate ourselves to move forward in partnership with Indigenous communities in a spirit of Reconciliation and collaboration.};  
      % \node[align=center, anchor=south,
      % above=0.5cm of land_acknowledgement,
      % text=black,
      % font=\bfseries] {\insertdate};
  \end{tikzpicture}
  \setbeamercolor{title}{fg=title_page_title_colour}
  \setbeamerfont{title}{size=\Large,series=\bfseries,family=\usefont{T1}{phv}{b}{n}}
  \setbeamercolor{author}{fg=title_page_author_colour}
  \setbeamerfont{author}{size=\large,series=\bfseries,family=\usefont{T1}{phv}{b}{n}}
  \setbeamercolor{institute}{fg=title_page_institute_colour}
  \setbeamerfont{institute}{size=\large,series=\bfseries,family=\usefont{T1}{phv}{b}{n}}
  \setbeamercolor{date}{fg=title_page_date_colour}
  \setbeamerfont{date}{series=\bfseries,family=\usefont{T1}{phv}{b}{n}}
	\titlepage
\end{frame}
\addtocounter{page}{-1}
}
% The outline page, with figure
% \newcommand{\outlinepage}[1]{%
% \begin{frame}[noframenumbering,plain]
%   \begin{tikzpicture}[remember picture,overlay]
%     \node[above right,inner sep=0pt,opacity=0.2] at (current page.south west)
%     {
%         \includegraphics[height=\paperheight,width=\paperwidth]{#1}
%     };
%   \end{tikzpicture}
%   \setbeamercolor{section in toc}{fg=outline_page_list_colour}
%   \setbeamerfont{section in toc}{size=\Large,series=\bfseries,family=\sffamily}
%   \frametitle{\textcolor{outline_page_title_colour}{\LARGE\bfseries Outline}}
%   \tableofcontents[hideallsubsections]
% \end{frame}
% \addtocounter{page}{-1}
% }
% The outline page, with figure
\newcommand{\outlinepage}[1]{%
\begin{frame}[noframenumbering,plain]
  \begin{tikzpicture}[remember picture,overlay]
    \node[above right,inner sep=0pt,opacity=0.2] at (current page.south west)
    {
        \includegraphics[height=\paperheight,width=\paperwidth]{#1}
    };
  \end{tikzpicture}
  \setbeamercolor{section in toc}{fg=outline_page_list_colour}
  % Use Helvetica Bold only for the outline slide TOC
  \setbeamerfont{section in toc}{size=\Large,family=\usefont{T1}{phv}{b}{n}}
  % Use Helvetica Bold for the outline title
  \frametitle{\textcolor{outline_page_title_colour}{\usefont{T1}{phv}{b}{n}\LARGE Outline}}
  \tableofcontents[hideallsubsections]
\end{frame}
\addtocounter{page}{-1}
}


%\let\oldsection\section
%\renewcommand{\section}[2]{\oldsection[#1]\newSectionSlide[#2]}


%%%%%%%%%%%%%%%%%%%%%
% CUSTOM SLIDE BACKGROUNDS
%%%%%%%%%%%%%%%%%%%%%
% Define custom background templates for different colors
\defbeamertemplate*{background canvas}{blue}{%
  \color{blue!15}\vrule width\paperwidth height\paperheight%
}
\defbeamertemplate*{background canvas}{green}{%
  \color{green!15}\vrule width\paperwidth height\paperheight%
}
\defbeamertemplate*{background canvas}{red}{%
  \color{red!15}\vrule width\paperwidth height\paperheight%
}
\defbeamertemplate*{background canvas}{yellow}{%
  \color{yellow!20}\vrule width\paperwidth height\paperheight%
}
\defbeamertemplate*{background canvas}{purple}{%
  \color{purple!15}\vrule width\paperwidth height\paperheight%
}
\defbeamertemplate*{background canvas}{orange}{%
  \color{orange!20}\vrule width\paperwidth height\paperheight%
}

% Define keys for the different background options
\makeatletter
\define@key{beamerframe}{blue}[true]{\setbeamertemplate{background canvas}[blue]}
\define@key{beamerframe}{green}[true]{\setbeamertemplate{background canvas}[green]}
\define@key{beamerframe}{red}[true]{\setbeamertemplate{background canvas}[red]}
\define@key{beamerframe}{yellow}[true]{\setbeamertemplate{background canvas}[yellow]}
\define@key{beamerframe}{purple}[true]{\setbeamertemplate{background canvas}[purple]}
\define@key{beamerframe}{orange}[true]{\setbeamertemplate{background canvas}[orange]}
\makeatother

% Reset to normal background for all frames by default
\BeforeBeginEnvironment{frame}{\setbeamertemplate{background canvas}[default]}

\newcommand{\punnett}[2]{
    \begin{center}
    \renewcommand{\arraystretch}{1.5} % Add space to rows
    \begin{tabular}{|c c | c|c|}
        \multicolumn{2}{c}{} & \multicolumn{2}{c}{\textbf{Father}} \\
        \multicolumn{2}{c}{} & #1 \\ \hline
        #2
    \end{tabular}
    \end{center}
}

% Colour definitions for Punnett squares
\usepackage{multirow}
\usepackage{colortbl}
\definecolor{punnettorange}{HTML}{F8C471} % A soft orange
\definecolor{punnettblack}{HTML}{BFC9CA} % A light grey for black
\definecolor{punnetttortie}{HTML}{E59866} % A brownish orange for tortie


\usetikzlibrary{positioning, decorations.pathreplacing, arrows.meta}

% --- TIKZ STYLES FOR CAT DIAGRAMS ---
\tikzset{
    % This is the main placeholder style for the cat images
    catnode/.style={
        draw, 
        rectangle, 
        rounded corners, 
        minimum height=1.5cm, 
        minimum width=2.5cm, 
        align=center, 
        font=\small\bfseries
    },
    % Allele styles (the X and Y at the top/side)
    allele/.style={font=\Large\bfseries},
    female/.style={allele, color=purple!80!black},
    male/.style={allele, color=blue!80!black},
    % --- Placeholder styles for different cats ---
    % YOU CAN REPLACE THE CONTENTS OF THESE NODES WITH YOUR IMAGES
    orange_cat/.style={
        catnode, 
        fill=orange!30, 
        text=black
    },
    black_cat/.style={
        catnode, 
        fill=black!70, 
        text=white
    },
    tortie_cat/.style={
        catnode, 
        fill=orange!50!black, % A mix for tortoiseshell
        text=white
    }
}

\usetikzlibrary{
    arrows.meta, % For nicer arrow heads (e.g., -Latex)
    positioning, % For relative node placement (e.g., above=of)
    automata     % For state diagrams, loops
}
% --- END TIKZ STYLES ---


\usecolortheme{orchid}
%% Listings
\usepackage{listings}
\definecolor{mygreen}{rgb}{0,0.6,0}
\definecolor{mygray}{rgb}{0.5,0.5,0.5}
\definecolor{mymauve}{rgb}{0.58,0,0.82}
\definecolor{mygold}{rgb}{1,0.843,0}
\definecolor{myblue}{rgb}{0.537,0.812,0.941}

\definecolor{mygold2}{RGB}{120,105,22}
\definecolor{mygrey2}{RGB}{50,50,50}

\definecolor{lgreen}{rgb}{0.6,0.9,.6}
\definecolor{lred}{rgb}{1,0.5,.5}

\lstloadlanguages{R}
\lstset{ %
  language=R,
  backgroundcolor=\color{black!05},   % choose the background color
  basicstyle=\footnotesize\ttfamily,        % size of fonts used for the code
  breaklines=true,                 % automatic line breaking only at whitespace
  captionpos=b,                    % sets the caption-position to bottom
  commentstyle=\color{mygreen},    % comment style
  escapeinside={\%*}{*)},          % if you want to add LaTeX within your code
  keywordstyle=\color{red},       % keyword style
  stringstyle=\color{mygold},     % string literal style
  keepspaces=true,
  columns=fullflexible,
  tabsize=4,
}
% Could also do (in lstset)
% basicstyle==\fontfamily{pcr}\footnotesize
\lstdefinelanguage{Renhanced}%
  {keywords={abbreviate,abline,abs,acos,acosh,action,add1,add,%
      aggregate,alias,Alias,alist,all,anova,any,aov,aperm,append,apply,%
      approx,approxfun,apropos,Arg,args,array,arrows,as,asin,asinh,%
      atan,atan2,atanh,attach,attr,attributes,autoload,autoloader,ave,%
      axis,backsolve,barplot,basename,besselI,besselJ,besselK,besselY,%
      beta,binomial,body,box,boxplot,break,browser,bug,builtins,bxp,by,%
      c,C,call,Call,case,cat,category,cbind,ceiling,character,char,%
      charmatch,check,chol,chol2inv,choose,chull,class,close,cm,codes,%
      coef,coefficients,co,col,colnames,colors,colours,commandArgs,%
      comment,complete,complex,conflicts,Conj,contents,contour,%
      contrasts,contr,control,helmert,contrib,convolve,cooks,coords,%
      distance,coplot,cor,cos,cosh,count,fields,cov,covratio,wt,CRAN,%
      create,crossprod,cummax,cummin,cumprod,cumsum,curve,cut,cycle,D,%
      data,dataentry,date,dbeta,dbinom,dcauchy,dchisq,de,debug,%
      debugger,Defunct,default,delay,delete,deltat,demo,de,density,%
      deparse,dependencies,Deprecated,deriv,description,detach,%
      dev2bitmap,dev,cur,deviance,off,prev,,dexp,df,dfbetas,dffits,%
      dgamma,dgeom,dget,dhyper,diag,diff,digamma,dim,dimnames,dir,%
      dirname,dlnorm,dlogis,dnbinom,dnchisq,dnorm,do,dotplot,double,%
      download,dpois,dput,drop,drop1,dsignrank,dt,dummy,dump,dunif,%
      duplicated,dweibull,dwilcox,dyn,edit,eff,effects,eigen,else,%
      emacs,end,environment,env,erase,eval,equal,evalq,example,exists,%
      exit,exp,expand,expression,External,extract,extractAIC,factor,%
      fail,family,fft,file,filled,find,fitted,fivenum,fix,floor,for,%
      For,formals,format,formatC,formula,Fortran,forwardsolve,frame,%
      frequency,ftable,ftable2table,function,gamma,Gamma,gammaCody,%
      gaussian,gc,gcinfo,gctorture,get,getenv,geterrmessage,getOption,%
      getwd,gl,glm,globalenv,gnome,GNOME,graphics,gray,grep,grey,grid,%
      gsub,hasTsp,hat,heat,help,hist,home,hsv,httpclient,I,identify,if,%
      ifelse,Im,image,\%in\%,index,influence,measures,inherits,install,%
      installed,integer,interaction,interactive,Internal,intersect,%
      inverse,invisible,IQR,is,jitter,kappa,kronecker,labels,lapply,%
      layout,lbeta,lchoose,lcm,legend,length,levels,lgamma,library,%
      licence,license,lines,list,lm,load,local,locator,log,log10,log1p,%
      log2,logical,loglin,lower,lowess,ls,lsfit,lsf,ls,machine,Machine,%
      mad,mahalanobis,make,link,margin,match,Math,matlines,mat,matplot,%
      matpoints,matrix,max,mean,median,memory,menu,merge,methods,min,%
      missing,Mod,mode,model,response,mosaicplot,mtext,mvfft,na,nan,%
      names,omit,nargs,nchar,ncol,NCOL,new,next,NextMethod,nextn,%
      nlevels,nlm,noquote,NotYetImplemented,NotYetUsed,nrow,NROW,null,%
      numeric,\%o\%,objects,offset,old,on,Ops,optim,optimise,optimize,%
      options,or,order,ordered,outer,package,packages,page,pairlist,%
      pairs,palette,panel,par,parent,parse,paste,path,pbeta,pbinom,%
      pcauchy,pchisq,pentagamma,persp,pexp,pf,pgamma,pgeom,phyper,pico,%
      pictex,piechart,Platform,plnorm,plogis,plot,pmatch,pmax,pmin,%
      pnbinom,pnchisq,pnorm,points,poisson,poly,polygon,polyroot,pos,%
      postscript,power,ppoints,ppois,predict,preplot,pretty,Primitive,%
      print,prmatrix,proc,prod,profile,proj,prompt,prop,provide,%
      psignrank,ps,pt,ptukey,punif,pweibull,pwilcox,q,qbeta,qbinom,%
      qcauchy,qchisq,qexp,qf,qgamma,qgeom,qhyper,qlnorm,qlogis,qnbinom,%
      qnchisq,qnorm,qpois,qqline,qqnorm,qqplot,qr,Q,qty,qy,qsignrank,%
      qt,qtukey,quantile,quasi,quit,qunif,quote,qweibull,qwilcox,%
      rainbow,range,rank,rbeta,rbind,rbinom,rcauchy,rchisq,Re,read,csv,%
      csv2,fwf,readline,socket,real,Recall,rect,reformulate,regexpr,%
      relevel,remove,rep,repeat,replace,replications,report,require,%
      resid,residuals,restart,return,rev,rexp,rf,rgamma,rgb,rgeom,R,%
      rhyper,rle,rlnorm,rlogis,rm,rnbinom,RNGkind,rnorm,round,row,%
      rownames,rowsum,rpois,rsignrank,rstandard,rstudent,rt,rug,runif,%
      rweibull,rwilcox,sample,sapply,save,scale,scan,scan,screen,sd,se,%
      search,searchpaths,segments,seq,sequence,setdiff,setequal,set,%
      setwd,show,sign,signif,sin,single,sinh,sink,solve,sort,source,%
      spline,splinefun,split,sqrt,stars,start,stat,stem,step,stop,%
      storage,strstrheight,stripplot,strsplit,structure,strwidth,sub,%
      subset,substitute,substr,substring,sum,summary,sunflowerplot,svd,%
      sweep,switch,symbol,symbols,symnum,sys,status,system,t,table,%
      tabulate,tan,tanh,tapply,tempfile,terms,terrain,tetragamma,text,%
      time,title,topo,trace,traceback,transform,tri,trigamma,trunc,try,%
      ts,tsp,typeof,unclass,undebug,undoc,union,unique,uniroot,unix,%
      unlink,unlist,unname,untrace,update,upper,url,UseMethod,var,%
      variable,vector,Version,vi,warning,warnings,weighted,weights,%
      which,while,window,write,\%x\%,x11,X11,xedit,xemacs,xinch,xor,%
      xpdrows,xy,xyinch,yinch,zapsmall,zip},%
   otherkeywords={!,!=,~,$,*,\%,\&,\%/\%,\%*\%,\%\%,<-,<<-,_,/},%
   alsoother={._$},%
   sensitive,%
   morecomment=[l]\#,%
   morestring=[d]",%
   morestring=[d]'% 2001 Robert Denham
  }%

%%%%%%% 
%% Definitions in yellow boxes
\usepackage{etoolbox}
\setbeamercolor{block title}{use=structure,fg=structure.fg,bg=structure.fg!40!bg}
\setbeamercolor{block body}{parent=normal text,use=block title,bg=block title.bg!20!bg}

\BeforeBeginEnvironment{definition}{%
	\setbeamercolor{block title}{fg=black,bg=yellow!20!white}
	\setbeamercolor{block body}{fg=black, bg=yellow!05!white}
}
\AfterEndEnvironment{definition}{
	\setbeamercolor{block title}{use=structure,fg=structure.fg,bg=structure.fg!20!bg}
	\setbeamercolor{block body}{parent=normal text,use=block title,bg=block title.bg!50!bg, fg=black}
}
\BeforeBeginEnvironment{importanttheorem}{%
	\setbeamercolor{block title}{fg=black,bg=red!20!white}
	\setbeamercolor{block body}{fg=black, bg=red!05!white}
}
\AfterEndEnvironment{importanttheorem}{
	\setbeamercolor{block title}{use=structure,fg=structure.fg,bg=structure.fg!20!bg}
	\setbeamercolor{block body}{parent=normal text,use=block title,bg=block title.bg!50!bg, fg=black}
}
\BeforeBeginEnvironment{importantproperty}{%
	\setbeamercolor{block title}{fg=black,bg=red!50!white}
	\setbeamercolor{block body}{fg=black, bg=red!30!white}
}
\AfterEndEnvironment{importantproperty}{
	\setbeamercolor{block title}{use=structure,fg=structure.fg,bg=structure.fg!20!bg}
	\setbeamercolor{block body}{parent=normal text,use=block title,bg=block title.bg!50!bg, fg=black}
}

% Colour for the outline page
\definecolor{outline_colour}{RGB}{230,165,83}
%% Colours for sections, subsections aand subsubsections
\definecolor{section_colour}{RGB}{27,46,28}
\definecolor{subsection_colour}{RGB}{52,128,56}
\definecolor{subsubsection_colour}{RGB}{150,224,154}
\definecolor{subsub_header_section}{RGB}{196,44,27}
%\definecolor{mygold}{rgb}{1,0.843,0}
% Beginning of a section
% \AtBeginSection[]{
% 	{
% 	  \setbeamercolor{section in toc}{fg=mygold}
% 		\setbeamercolor{background canvas}{bg=section_colour}
% 		\begin{frame}[noframenumbering,plain]
% 			\framesubtitle{\nameofthepart Chapter \insertromanpartnumber \ -- \iteminsert{\insertpart}}
% 			\tableofcontents[
% 				currentsection,
% 				sectionstyle=show/shaded,
% 				subsectionstyle=show/hide/hide,
% 				subsubsectionstyle=hide/hide/hide]
% 		\end{frame}
% 	\addtocounter{page}{-1}
% 	%\addtocounter{framenumber}{-1} 
% 	}
% }


% % Beginning of a section
% \AtBeginSubsection[]{
% 	{
% 	  \setbeamercolor{section in toc}{fg=mygold}
% 		\setbeamercolor{background canvas}{bg=subsection_colour}
% 		\begin{frame}[noframenumbering,plain]
% 				\framesubtitle{\nameofthepart Chapter \insertromanpartnumber \ -- \iteminsert{\insertpart}}
% 				\tableofcontents[
% 					currentsection,
% 					sectionstyle=show/hide,
% 					currentsubsection,
% 					subsectionstyle=show/shaded/hide,
% 					subsubsectionstyle=show/hide/hide]
% 			\end{frame}
% 		\addtocounter{page}{-1}
% 	}
% }

% \newcommand{\newSubSectionSlide}[1]{
% \begin{frame}[noframenumbering,plain]
%   \begin{tikzpicture}[remember picture,overlay]
%     \node[above right,inner sep=0pt,opacity=0.2] at (current page.south west)
%     {
%         \includegraphics[height=\paperheight,width=\paperwidth]{#1}
%     };
%   \end{tikzpicture}
%   \setbeamercolor{section in toc}{fg=subsub_header_section}
%   \setbeamerfont{section in toc}{size=\Large,series=\bfseries}
%   \setbeamertemplate{section in toc shaded}[default][60]
%   \setbeamertemplate{subsection in toc shaded}[default][60]
%   %\setbeamercolor{background canvas}{bg=section_colour}
%   \tableofcontents[
%     currentsection,
%     sectionstyle=show/hide,
%     currentsubsection,
%     subsectionstyle=show/shaded/hide,
%     subsubsectionstyle=show/hide/hide]
% \end{frame}
% \addtocounter{page}{-1}
% }


% % Beginning of a section
% \AtBeginSubsubsection[]{
% 	{
% 	  \setbeamercolor{section in toc}{fg=subsub_header_section}
% 	  \setbeamercolor{subsubsection in toc}{fg=mygold2}
% 	  \setbeamercolor{subsubsection in toc shaded}{fg=mygrey2}
% 		\setbeamercolor{background canvas}{bg=subsubsection_colour}
% 		\begin{frame}[noframenumbering,plain]
% 				\framesubtitle{\nameofthepart Chapter \insertromanpartnumber \ -- \iteminsert{\insertpart}}
% 				\tableofcontents[
% 					currentsection,
% 					sectionstyle=show/hide,
% 					currentsubsection,
% 					subsectionstyle=show/hide/shaded
% 					currentsubsubsection]%,
% 					%subsubsectionstyle=hide/hide/shaded]
% 					%currentsubsubsection]
% 			\end{frame}
% 		\addtocounter{page}{-1}
% 	}
% }


\title{Clustering \& Classification using ANNs}
\author{\texorpdfstring{Julien Arino\newline University of Manitoba\newline\url{julien.arino@umanitoba.ca}}{Julien Arino}}
\date{}
\IfFileExists{upquote.sty}{\usepackage{upquote}}{}
\begin{document}

%%%%%%%%%%%%%%%%%%%%%%%%%%%%%%%%%
%%%%%%%%%%%%%%%%%%%%%%%%%%%%%%%%%
%% TITLE AND OUTLINE
%%%%%%%%%%%%%%%%%%%%%%%%%%%%%%%%%
%%%%%%%%%%%%%%%%%%%%%%%%%%%%%%%%%
\titlepagewithfigure{FIGS-transitions/Gemini_Generated_Image_pbzl1mpbzl1mpbzl.jpeg}
\outlinepage{FIGS-transitions/Gemini_Generated_Image_yegaxcyegaxcyega.jpeg}

%%%%%%%%%%%%%%%%%%%
%%%%%%%%%%%%%%%%%%%
%%%%%%%%%%%%%%%%%%%
%%%%%%%%%%%%%%%%%%%
%%%%%%%%%%%%%%%%%%%
%%%%%%%%%%%%%%%%%%%
%%%%%%%%%%%%%%%%%%%
%%%%%%%%%%%%%%%%%%%
\section{Neural networks (the perceptron)}
\newSectionSlide{FIGS-transitions/Gemini_Generated_Image_5g0m5e5g0m5e5g0m.jpeg}

\begin{frame}{Artificial neural network (ANN) - from \href{https://en.wikipedia.org/wiki/Artificial_neural_network}{Wikipedia}}
    \begin{quote}
    \textbf{Artificial neural networks} (ANNs) are computing systems inspired by the biological neural networks that constitute animal brains
    \vfill
    An ANN is based on a collection of connected units or nodes called \textbf{artificial neurons}, which loosely model the neurons in a biological brain. Each connection, like the synapses in a biological brain, can transmit a signal to other neurons. An artificial neuron receives signals then processes them and can signal neurons connected to it. The ``signal" at a connection is a real number, and the output of each neuron is computed by some non-linear function of the sum of its inputs. The connections are called edges. Neurons and edges typically have a weight that adjusts as learning proceeds. The weight increases or decreases the strength of the signal at a connection. Neurons may have a threshold such that a signal is sent only if the aggregate signal crosses that threshold.
    \end{quote}
\end{frame}


\begin{frame}{The perceptron}
    One of the first neural networks (invented 1943, implemented 1957), made for simple classification tasks, for example recognising letters or numbers
    \vfill
    Two layers: the \emph{input} layer (the \emph{retina}) and the \emph{output} layer
    \vfill
    Inputs are 0 or 1, so are outputs
\end{frame}


\begin{frame}{}
    \begin{center}
        \def\hhskip{8cm}
        \def\vvskip{2.75cm}
        \begin{tikzpicture}[scale=0.75, 
		every node/.style={transform shape},
		auto,
		cloud/.style={minimum width={width("N-1")+2pt},
			draw, ellipse},
		connected/.style={dotted,-}]
		%% Input nodes
		\node [cloud] at (0,0*\vvskip) (i1) {$0\lor 1$};
		\node [cloud] at (0,-1*\vvskip) (i2) {$0\lor 1$};
		\node [cloud] at (0,-2*\vvskip) (i3) {$0\lor 1$};
		\node [cloud] at (0,-3*\vvskip) (i4) {$0\lor 1$};
		%% Output nodes
		\node [cloud] at (1*\hhskip,-0.5*\vvskip) (o1) {$\sum$};
		\node [cloud] at (1*\hhskip,-1.5*\vvskip) (o2) {$\sum$};
		\node [cloud] at (1*\hhskip,-2.5*\vvskip) (o3) {$\sum$};
		%% Arcs
		\path [line] (i1) to node [midway, above] (TextNode) {} (o1);
		\path [line] (i1) to node [midway, above] (TextNode) {} (o2);
		\path [line] (i1) to node [midway, above] (TextNode) {} (o3);
		\path [line] (i2) to node [midway, above] (TextNode) {} (o1);
		\path [line] (i2) to node [midway, above] (TextNode) {} (o2);
		\path [line] (i2) to node [midway, above] (TextNode) {} (o3);
		\path [line] (i3) to node [midway, above] (TextNode) {} (o1);
		\path [line] (i3) to node [midway, above] (TextNode) {} (o2);
		\path [line] (i3) to node [midway, above] (TextNode) {} (o3);
		\path [line] (i4) to node [midway, above] (TextNode) {} (o1);
		\path [line] (i4) to node [midway, above] (TextNode) {} (o2);
		\path [line] (i4) to node [midway, above] (TextNode) {} (o3);
        %% Arrows out of output nodes
		\path [line] (o1) to node [midway, above] (TextNode) {$0\lor 1$} (1.2*\hhskip,-0.5*\vvskip);
		\path [line] (o2) to node [midway, above] (TextNode) {$0\lor 1$} (1.2*\hhskip,-1.5*\vvskip);
		\path [line] (o3) to node [midway, above] (TextNode) {$0\lor 1$} (1.2*\hhskip,-2.5*\vvskip);
        %% Some labels
        \node[align=left] at (0.5*\hhskip,0) {Connections};
        \node[align=left] at (0*\hhskip,0.33*\vvskip) {Retina};
        \node[align=left] at (1*\hhskip,0.33*\vvskip) {Output layer};
	\end{tikzpicture}
    \end{center}
    The connections into the output layer are called \emph{synapses}, they are modifiable
\end{frame}

\begin{frame}{The activation function}
    \begin{center}
        \def\hhskip{8cm}
        \def\vvskip{2cm}
        \begin{tikzpicture}[scale=0.75, 
		every node/.style={transform shape},
		auto,
		cloud/.style={minimum width={width("N-1")+2pt},
			draw, ellipse},
		connected/.style={dotted,-}]
		%% Input nodes
		\node [cloud] at (0,0*\vvskip) (i1) {$x_0$};
		\node [cloud] at (0,-1*\vvskip) (i2) {$x_1$};
		\node [cloud] at (0,-2*\vvskip) (i3) {$x_i$};
		\node [cloud] at (0,-3*\vvskip) (i4) {$x_\ell$};
		%% Output nodes
		\node [cloud] at (1*\hhskip,-1.5*\vvskip) (o1) {$a_j=\sum_{i=1}^I w_{ij}x_i$};
		%% Arcs
		\path [line] (i1) to node [midway, above] (TextNode) {$w_{1j}$} (o1);
		\path [line] (i2) to node [midway, above] (TextNode) {$w_{2j}$} (o1);
		\path [line] (i3) to node [midway, above] (TextNode) {$w_{ij}$} (o1);
		\path [line] (i4) to node [midway, above] (TextNode) {$w_{I j}$} (o1);
        %% Arrows out of output nodes
		\path [line] (o1) to node [midway, above] (TextNode) {$o_j=f(a_j)$} (1.5*\hhskip,-1.5*\vvskip);
	\end{tikzpicture}
    \end{center}
    Here,
    \[
        f(a_j)= \begin{cases}
            0 & \text{if }a_j\leq 0 \\
            1 & \text{if }a_j>0
        \end{cases}
    \]
\end{frame}

\begin{frame}{The activation function}
    We have $I$ input neurons taking values $0$ or $1$, $O$ output neurons taking values $0$ or $1$, weights $W=[w_{ij}]\in\M_{IO}$ and a threshold function $f$
    \vfill
    More generally, use a threshold $\theta_j$ for each output neuron
    \[
        o_j= \begin{cases}
            0 & \text{if }a_j\leq \theta_j \\
            1 & \text{if }a_j>\theta_j
        \end{cases}
    \]
    \vfill
    The thresholds (or \emph{response bias}) and the weights are modifiable by learning. To do that easily for the threshold, consider an input neuron that is always on, say neuron 0, and set weights $w_{0j}=-\theta_j$, making the weights matrix an $(I+1)\times O$-matrix
\end{frame}


\begin{frame}
    Another way to write the activation is 
    \[
        o_j= \begin{cases}
            0 & \text{if }a_j+w_{0j}\leq 0 \\
            1 & \text{if }a_j+w_{0j}>0
        \end{cases}
    \]
    where $w_{0j}=-\theta_j$
\end{frame}


\begin{frame}{Learning something simple}
    The aim is to adjust the synaptic weights so that the proper response is provided to a given stimulus
    \vfill
    Let us first do a simple example: the OR truth table
    \begin{center}
        \begin{tabular}{cccc}
            0 & 0 & $\mapsto$ & 0 \\
            1 & 0 & $\mapsto$ & 1 \\
            0 & 1 & $\mapsto$ & 1 \\
            1 & 1 & $\mapsto$ & 1 \\
        \end{tabular}
    \end{center}
    So we have two neurons in the retina and a single output neuron
\end{frame}

\begin{frame}{Supervised learning}
    (From R. Rojas)
    \begin{quote}
    \textbf{Supervised learning}: method in which some input vectors are collected and presented to the network. The output computed by the network is observed and the deviation from the expected answer is measured
    \vskip0.5cm
    The weights are corrected according to the magnitude of the error in the way defined by the learning algorithm
    \vskip0.5cm
    Also called \emph{learning with a teacher}, since a control process knows the correct answer for the set of selected input vectors
    \end{quote}
\end{frame}

\begin{frame}{Further distinctions in supervised learning methods}
    Methods with \emph{reinforcement} or \emph{error correction} 
    \vfill
    \begin{itemize}
        \item \textbf{Reinforcement learning}: used when after each presentation of an input-output example, we only know whether the network produces the desired result or not. 
        Weights are updated based on this information (i.e., the Boolean values true or false), so only the input vector can be used for weight correction
        \vfill
        \item In learning with \textbf{error correction}, the
        magnitude of the error, together with the input vector, determines the magnitude of the corrections to the weights. In many cases, we try to eliminate the error in a single correction step
    \end{itemize}
\end{frame}

\begin{frame}{A first learning algorithm}
    Suppose the training set consists of two sets of points $P$ and $N$
    \vfill
    \begin{itemize}
        \item \textbf{start:} Generate random weight vector $w_0$; set $t := 0$
        \vfill
        \item \textbf{test:} A vector $x\in P\cup N$ is selected randomly
        \begin{itemize}
            \item if $x\in P$ and $\langle w_t,x\rangle > 0$ go to \textbf{test}
            \item if $x\in P$ and $\langle w_t,x\rangle\leq 0$ go to \textbf{add}
            \item if $x\in N$ and $\langle w_t,x\rangle < 0$ go to \textbf{test}
            \item if $x\in N$ and $\langle w_t,x\rangle \geq 0$ go to \textbf{subtract}
        \end{itemize}
        \vfill
        \item \textbf{add:} set $w_{t+1} = w_t + x$ and $t := t + 1$, goto \textbf{test}
        \vfill
        \item \textbf{subtract:} set $w_{t+1} = w_t - x$ and $t := t + 1$, goto \textbf{test}
    \end{itemize}
\end{frame}

\begin{frame}{Widrow-Hoff learning rule}
    Need to provide the correct answer, i.e., this is a supervised learning rule
    \vfill
    An output cell only learns if it is mistaken
    \vfill
    Present random inputs and apply the rule if the output does not match the known output
\end{frame}

\begin{frame}{Widrow-Hoff learning rule}
    \begin{equation}\label{eq:WidrowHoff}
        w_{ij}^{(t+1)} = w_{ij}^{(t)}+\eta(t_j-o_j)x_j = w_{ij}^{(t)}+\Delta w_{ij}        
    \end{equation}
    with
    \begin{itemize}
        \item $\Delta w_{ij}$ correction to add to the weight $w_{ij}$
        \item $x_i$: value ($0$ or $1$) of the $i$th retinal cell
        \item $o_j$: response of the $j$th output cell
        \item $t_j$ target response (correct desired response)
        \item $w_{ij}^{(t)}$: weight of the synapse between the $i$th retinal cell and $j$th output cell at time $t$. Typically initiated at small random values
        \item $\eta$: small positive constant, the \emph{learning constant}
    \end{itemize} 
\end{frame}

\begin{frame}{Learning OR}
    \begin{center}
        \begin{tabular}{cccc}
            0 & 0 & $\mapsto$ & 0 \\
            1 & 0 & $\mapsto$ & 1 \\
            0 & 1 & $\mapsto$ & 1 \\
            1 & 1 & $\mapsto$ & 1 \\
        \end{tabular}
    \end{center}
    Three cells in the retina (two inputs and the ``dummy'' cell used for the threshold) and one output cell. So inputs and outputs must be
    \begin{center}
        \begin{tabular}{ccccc}
            1 & 0 & 0 & $\mapsto$ & 0 \\
            1 & 1 & 0 & $\mapsto$ & 1 \\
            1 & 0 & 1 & $\mapsto$ & 1 \\
            1 & 1 & 1 & $\mapsto$ & 1 \\
        \end{tabular}
    \end{center}
    Initialise the $3\times 1$ weight matrix $W$ to zero:
    \[
        W=\begin{pmatrix}
            w_0 = -\theta \\
            w_1 \\ w_2
        \end{pmatrix}
        = \begin{pmatrix}
            0\\ 0\\ 0
        \end{pmatrix}
    \]
\end{frame}


\begin{frame}{Procedure}
    We choose one random association in 
    \begin{center}
        \begin{tabular}{ccccc}
            1 & 0 & 0 & $\mapsto$ & 0 \\
            1 & 1 & 0 & $\mapsto$ & 1 \\
            1 & 0 & 1 & $\mapsto$ & 1 \\
            1 & 1 & 1 & $\mapsto$ & 1 \\
        \end{tabular}
    \end{center}
    say, the fourth one. So we present $[1,1,1]$ and expect an output of 1. We have
    \[
        a = \sum_i w_ix_i = (1\times 0)+(1\times 0)+(1\times 0) = 0
    \]
    This being $\leq 0$ means that $o=0$, giving an error of 1
\end{frame}

\begin{frame}{Applying the rule}
    Suppose the learning constant $\eta=0.1$. Then applying \eqref{eq:WidrowHoff},
    \begin{align*}
        \Delta w_0 &= \eta(t-o)x_0 = 0.1\times(1-0)\times 1= 0.1 \\
        \Delta w_1 &= \eta(t-o)x_0 = 0.1\times(1-0)\times 1= 0.1 \\
        \Delta w_2 &= \eta(t-o)x_0 = 0.1\times(1-0)\times 1= 0.1         
    \end{align*}
    Applying the correction, $W$ becomes
    \[
        W=\begin{pmatrix}
            0.1\\ 0.1\\ 0.1
        \end{pmatrix}
    \]
\end{frame}

\begin{frame}{Trying another input}
    Suppose we now present the first input $[1,0,0]$. This should produce a result of 0. Then
    \[
        a = \sum_i w_ix_i = (1\times 0.1)+(0\times 0.1)+(0\times 0.1) = 0.1
    \]
    which is $>0$, so $o=1$. We compute the correction
    \begin{align*}
        \Delta w_0 &= \eta(t-o)x_0 = 0.1\times(0-1)\times 1= -0.1 \\
        \Delta w_1 &= \eta(t-o)x_1 = 0.1\times(0-1)\times 0= 0 \\
        \Delta w_1 &= \eta(t-o)x_2 = 0.1\times(0-1)\times 0= 0         
    \end{align*}
    and adjust the weights, giving
    \[
        W=\begin{pmatrix}
            0\\ 0.1\\ 0.1
        \end{pmatrix}
    \]
\end{frame}

\begin{frame}{And we are done!}
    With the weights
    \[
        W=\begin{pmatrix}
            0\\ 0.1\\ 0.1
        \end{pmatrix}
    \]
    we are done. Indeed
    \begin{center}
        \begin{tabular}{cccccc}
            Input 0 & Input 1 & Input 2 & a & o & Should be \\
            1 & 0 & 0 & 0 & 0 & 0 \\
            1 & 1 & 0 & 0+0.1+0 & 1 & 1 \\
            1 & 0 & 1 & 0+0+0.1 & 1 & 1 \\
            1 & 1 & 1 & 0+0.1+0.1 & 1 & 1 \\
        \end{tabular}
    \end{center}    
\end{frame}


\begin{frame}{Learning XOR}
    Let us now look at the XOR truth table
    \begin{center}
        \begin{tabular}{cccc}
            0 & 0 & $\mapsto$ & 0 \\
            1 & 0 & $\mapsto$ & 1 \\
            0 & 1 & $\mapsto$ & 1 \\
            1 & 1 & $\mapsto$ & 0 \\
        \end{tabular}
    \end{center}
    This problem is not solvable with a simple perceptron of the type we just used, as truth table is not \emph{linearly separable}
    \vfill
    Indeed, we would get weights $w_1>0$, $w_2>0$ to activate when presenting $[1,0]$ and $[0,1]$, but would require that the sum of the weights when applied to the input $[1,1]$, give a negative value.
\end{frame}

\begin{frame}{Linear separability and OR and XOR}
    \begin{center}
    \begin{tikzpicture}[scale=3,auto]
        \draw[step=1cm,gray!25!,very thin] (-0.25,-0.25) grid (1.25,1.25);
        \draw[thick,->] (-0.25,0) -- (1.5,0) node[anchor=north west] {$x_1$};
        \draw[thick,->] (0,-0.25) -- (0,1.5) node[anchor=south east] {$x_2$};
        \node[circle,draw,very thick] (c) at (0,0){};
        \fill (1,0)  circle[radius=2pt];
        \fill (0,1)  circle[radius=2pt];
        \fill (1,1)  circle[radius=2pt];
    \end{tikzpicture}\quad
    \begin{tikzpicture}[scale=3,auto]
        \draw[step=1cm,gray!25!,very thin] (-0.25,-0.25) grid (1.25,1.25);
        \draw[thick,->] (-0.25,0) -- (1.5,0) node[anchor=north west] {$x_1$};
        \draw[thick,->] (0,-0.25) -- (0,1.5) node[anchor=south east] {$x_2$};
        \node[circle,draw,very thick] (c) at (0,0){};
        \fill (1,0)  circle[radius=2pt];
        \fill (0,1)  circle[radius=2pt];
        \node[circle,draw,very thick] (c) at (1,1){};
    \end{tikzpicture}
    \end{center}
    A single-layer perceptron can only learn linearly separable problems
\end{frame}


\begin{frame}{Adding a hidden layer}
    It is possible to do XOR, but we need to add a \textbf{hidden layer}
    \begin{center}
        \def\hhskip{4cm}
        \def\vvskip{4cm}
        \begin{tikzpicture}[scale=1, 
		every node/.style={transform shape},
		auto,
		cloud/.style={minimum width={width("N-1")+2pt},
			draw, ellipse},
		connected/.style={dotted,-}]
		%% Input nodes
		\node [cloud] at (0,0*\vvskip) (i1) {$0\lor 1$};
		\node [cloud] at (0,-1*\vvskip) (i2) {$0\lor 1$};
        %% Hidden node
        \node [cloud] at (1*\hhskip,-0.5*\vvskip) (h1) {$\theta=1$};
		%% Output node
		\node [cloud] at (2*\hhskip,-0.5*\vvskip) (o1) {$\theta=0$};
		%% Arcs
		\path [line] (i1) to node [midway, above, sloped] (TextNode) {$w=1$} (o1);
		\path [line] (i2) to node [midway, below, sloped] (TextNode) {$w=1$} (o1);
		\path [line] (i1) to node [midway, below, sloped] (TextNode) {$w=0.6$} (h1);
		\path [line] (i2) to node [midway, above, sloped] (TextNode) {$w=0.6$} (h1);
		\path [line] (h1) to node [pos=0.25, above] (TextNode) {$w=-2$} (o1);
	\end{tikzpicture}
    \end{center}
\end{frame}

\begin{frame}[fragile]\frametitle{Using \code{neuralnet} to learn OR}
First, create the truth table
\begin{knitrout}
\definecolor{shadecolor}{rgb}{0.969, 0.969, 0.969}\color{fgcolor}\begin{kframe}
\begin{alltt}
\hldef{OR_table} \hlkwb{=} \hlkwd{matrix}\hldef{(}\hlkwd{c}\hldef{(}\hlnum{0}\hldef{,} \hlnum{0}\hldef{,} \hlnum{0}\hldef{,}
                    \hlnum{1}\hldef{,} \hlnum{0}\hldef{,} \hlnum{1}\hldef{,}
                    \hlnum{0}\hldef{,} \hlnum{1}\hldef{,} \hlnum{1}\hldef{,}
                    \hlnum{1}\hldef{,} \hlnum{1}\hldef{,} \hlnum{1}\hldef{),}
                  \hlkwc{nc} \hldef{=} \hlnum{3}\hldef{,} \hlkwc{byrow} \hldef{=} \hlnum{TRUE}\hldef{)}
\hldef{OR_table} \hlkwb{=} \hlkwd{as.data.frame}\hldef{(OR_table)}
\hlkwd{colnames}\hldef{(OR_table)} \hlkwb{=} \hlkwd{c}\hldef{(}\hlsng{"x1"}\hldef{,} \hlsng{"x2"}\hldef{,} \hlsng{"OR"}\hldef{)}
\end{alltt}
\end{kframe}
\end{knitrout}
\end{frame}

\begin{frame}[fragile]\frametitle{Now create and train the NN}
The ``formula'' is to find the OR column using the x1 and x2 columns. We use no hidden layer
\vfill
\begin{knitrout}
\definecolor{shadecolor}{rgb}{0.969, 0.969, 0.969}\color{fgcolor}\begin{kframe}
\begin{alltt}
\hldef{nn_OR} \hlkwb{=} \hlkwd{neuralnet}\hldef{(OR} \hlopt{~} \hldef{x1} \hlopt{+} \hldef{x2,}
                  \hlkwc{data} \hldef{= OR_table,}
                  \hlkwc{act.fct} \hldef{=} \hlsng{"logistic"}\hldef{,}
                  \hlkwc{hidden} \hldef{=} \hlnum{0}\hldef{,}
                  \hlkwc{linear.output} \hldef{=} \hlnum{FALSE}\hldef{)}
\hlcom{# Plot the result}
\hlkwd{plot}\hldef{(nn_OR,} \hlkwc{rep} \hldef{=} \hlsng{"best"}\hldef{)}
\end{alltt}
\end{kframe}
\end{knitrout}
\end{frame}

\maxFrameImage{FIGS/ann-plot-NN-for-OR-1.pdf}

\begin{frame}[fragile]\frametitle{Testing the result}
\begin{knitrout}
\definecolor{shadecolor}{rgb}{0.969, 0.969, 0.969}\color{fgcolor}\begin{kframe}
\begin{alltt}
\hldef{pred} \hlkwb{=} \hlkwd{predict}\hldef{(nn_OR, OR_table)}
\hldef{OR_table}\hlopt{$}\hldef{result} \hlkwb{=} \hldef{pred} \hlopt{>} \hlnum{0.5}
\hlkwd{kable}\hldef{(OR_table,} \hlsng{"latex"}\hldef{,} \hlkwc{booktabs} \hldef{=} \hlnum{TRUE}\hldef{)}
\end{alltt}
\end{kframe}
\begin{tabular}{rrrl}
\toprule
x1 & x2 & OR & result\\
\midrule
0 & 0 & 0 & FALSE\\
1 & 0 & 1 & TRUE\\
0 & 1 & 1 & TRUE\\
1 & 1 & 1 & TRUE\\
\bottomrule
\end{tabular}

\end{knitrout}
\end{frame}


\begin{frame}[fragile]\frametitle{Now the XOR truth table}
\begin{knitrout}
\definecolor{shadecolor}{rgb}{0.969, 0.969, 0.969}\color{fgcolor}\begin{kframe}
\begin{alltt}
\hldef{XOR_table} \hlkwb{=} \hlkwd{matrix}\hldef{(}\hlkwd{c}\hldef{(}\hlnum{0}\hldef{,} \hlnum{0}\hldef{,} \hlnum{0}\hldef{,}
                     \hlnum{1}\hldef{,} \hlnum{0}\hldef{,} \hlnum{1}\hldef{,}
                     \hlnum{0}\hldef{,} \hlnum{1}\hldef{,} \hlnum{1}\hldef{,}
                     \hlnum{1}\hldef{,} \hlnum{1}\hldef{,} \hlnum{0}\hldef{),}
                  \hlkwc{nc} \hldef{=} \hlnum{3}\hldef{,} \hlkwc{byrow} \hldef{=} \hlnum{TRUE}\hldef{)}
\hldef{XOR_table} \hlkwb{=} \hlkwd{as.data.frame}\hldef{(XOR_table)}
\hlkwd{colnames}\hldef{(XOR_table)} \hlkwb{=} \hlkwd{c}\hldef{(}\hlsng{"x1"}\hldef{,} \hlsng{"x2"}\hldef{,} \hlsng{"XOR"}\hldef{)}
\end{alltt}
\end{kframe}
\end{knitrout}
\end{frame}

\begin{frame}[fragile]\frametitle{Try to learn it without a hidden layer}
\begin{knitrout}
\definecolor{shadecolor}{rgb}{0.969, 0.969, 0.969}\color{fgcolor}\begin{kframe}
\begin{alltt}
\hldef{nn_XOR} \hlkwb{=} \hlkwd{neuralnet}\hldef{(XOR} \hlopt{~} \hldef{x1} \hlopt{+} \hldef{x2,}
                   \hlkwc{data} \hldef{= XOR_table,}
                   \hlkwc{act.fct} \hldef{=} \hlsng{"logistic"}\hldef{,}
                   \hlkwc{hidden} \hldef{=} \hlnum{0}\hldef{,}
                   \hlkwc{linear.output} \hldef{=} \hlnum{FALSE}\hldef{)}
\hldef{pred} \hlkwb{=} \hlkwd{predict}\hldef{(nn_XOR, XOR_table)}
\hldef{XOR_table}\hlopt{$}\hldef{result} \hlkwb{=} \hldef{pred} \hlopt{>} \hlnum{0.5}
\hlkwd{kable}\hldef{(XOR_table,} \hlsng{"latex"}\hldef{,} \hlkwc{booktabs} \hldef{=} \hlnum{TRUE}\hldef{)}
\end{alltt}
\end{kframe}
\begin{tabular}{rrrl}
\toprule
x1 & x2 & XOR & result\\
\midrule
0 & 0 & 0 & FALSE\\
1 & 0 & 1 & TRUE\\
0 & 1 & 1 & FALSE\\
1 & 1 & 0 & TRUE\\
\bottomrule
\end{tabular}

\end{knitrout}
\end{frame}

\begin{frame}[fragile]\frametitle{Now with a hidden layer}
\begin{knitrout}
\definecolor{shadecolor}{rgb}{0.969, 0.969, 0.969}\color{fgcolor}\begin{kframe}
\begin{alltt}
\hldef{nn_XOR} \hlkwb{=} \hlkwd{neuralnet}\hldef{(XOR} \hlopt{~} \hldef{x1} \hlopt{+} \hldef{x2,}
                   \hlkwc{data} \hldef{= XOR_table,}
                   \hlkwc{act.fct} \hldef{=} \hlsng{"tanh"}\hldef{,}
                   \hlkwc{hidden} \hldef{=} \hlnum{1}\hldef{)}
\hldef{pred} \hlkwb{=} \hlkwd{predict}\hldef{(nn_XOR, XOR_table)}
\hldef{XOR_table}\hlopt{$}\hldef{result} \hlkwb{=} \hldef{pred} \hlopt{>} \hlnum{0.5}
\hlkwd{kable}\hldef{(XOR_table,} \hlsng{"latex"}\hldef{,} \hlkwc{booktabs} \hldef{=} \hlnum{TRUE}\hldef{)}
\end{alltt}
\end{kframe}
\begin{tabular}{rrrl}
\toprule
x1 & x2 & XOR & result\\
\midrule
0 & 0 & 0 & FALSE\\
1 & 0 & 1 & TRUE\\
0 & 1 & 1 & TRUE\\
1 & 1 & 0 & TRUE\\
\bottomrule
\end{tabular}

\end{knitrout}
\vfill
Should look into options.. :)
\end{frame}

\begin{frame}[fragile]\frametitle{An example from the \code{neuralnet} manual -- Training vs testing sets}
\code{iris} is a built-in \code{R} dataset detailing physical characteristics of 150 flowers from 3 iris species
\vfill
\begin{knitrout}
\definecolor{shadecolor}{rgb}{0.969, 0.969, 0.969}\color{fgcolor}\begin{kframe}
\begin{alltt}
\hldef{train_idx} \hlkwb{<-} \hlkwd{sample}\hldef{(}\hlkwd{nrow}\hldef{(iris),} \hlnum{2}\hlopt{/}\hlnum{3} \hlopt{*} \hlkwd{nrow}\hldef{(iris))}
\hldef{iris_train} \hlkwb{<-} \hldef{iris[train_idx, ]}
\hldef{iris_test} \hlkwb{<-} \hldef{iris[}\hlopt{-}\hldef{train_idx, ]}
\end{alltt}
\end{kframe}
\end{knitrout}
\vfill
Thus we pick at random 2/3 of the data for training and 1/3 for testing
\end{frame}

\begin{frame}[fragile]\frametitle{An example from the \code{neuralnet} manual}
\begin{knitrout}
\definecolor{shadecolor}{rgb}{0.969, 0.969, 0.969}\color{fgcolor}\begin{kframe}
\begin{alltt}
\hldef{nn} \hlkwb{<-} \hlkwd{neuralnet}\hldef{(Species} \hlopt{==} \hlsng{"setosa"} \hlopt{~} \hldef{Petal.Length} \hlopt{+} \hldef{Petal.Width,}
                \hldef{iris_train,} \hlkwc{linear.output} \hldef{=} \hlnum{FALSE}\hldef{)}
\hldef{pred} \hlkwb{<-} \hlkwd{predict}\hldef{(nn, iris_test)}
\hlkwd{table}\hldef{(iris_test}\hlopt{$}\hldef{Species} \hlopt{==} \hlsng{"setosa"}\hldef{, pred[,} \hlnum{1}\hldef{]} \hlopt{>} \hlnum{0.5}\hldef{)}
\end{alltt}
\begin{verbatim}
##        
##         FALSE TRUE
##   FALSE    34    0
##   TRUE      0   16
\end{verbatim}
\end{kframe}
\end{knitrout}
\end{frame}

\begin{frame}[fragile]\frametitle{Another example -- multiclass classification}
\begin{knitrout}
\definecolor{shadecolor}{rgb}{0.969, 0.969, 0.969}\color{fgcolor}\begin{kframe}
\begin{alltt}
\hldef{nn} \hlkwb{<-} \hlkwd{neuralnet}\hldef{((Species} \hlopt{==} \hlsng{"setosa"}\hldef{)} \hlopt{+}
                  \hldef{(Species} \hlopt{==} \hlsng{"versicolor"}\hldef{)} \hlopt{+}
                  \hldef{(Species} \hlopt{==} \hlsng{"virginica"}\hldef{)}
                 \hlopt{~} \hldef{Petal.Length} \hlopt{+} \hldef{Petal.Width,}
                \hldef{iris_train,} \hlkwc{linear.output} \hldef{=} \hlnum{FALSE}\hldef{)}
\hldef{pred} \hlkwb{<-} \hlkwd{predict}\hldef{(nn, iris_test)}
\hlkwd{table}\hldef{(iris_test}\hlopt{$}\hldef{Species,} \hlkwd{apply}\hldef{(pred,} \hlnum{1}\hldef{, which.max))}
\end{alltt}
\begin{verbatim}
##             
##               1  2  3
##   setosa     16  0  0
##   versicolor  0 17  1
##   virginica   0  2 14
\end{verbatim}
\end{kframe}
\end{knitrout}
\end{frame}

\end{document}
